\section{Conclusions and Recommendations}
\label{sec:conclusions}
Viscereality increased reported spatial self-extension during paced breathing. The original predictions concerning reversed feedback, temporal experience and an objective shift in peripersonal space remain unresolved. Respiratory PPS data collection is complete and analysis remains ongoing.

\subsection{Expected and achieved outputs}
\label{sec:outputs}
% Required BIAL categories; application p.15 is the expected-count authority.
\begingroup
\small\setlength{\tabcolsep}{4pt}\renewcommand{\arraystretch}{1.05}
\begin{tabular}{p{.51\linewidth}r r}
\toprule
Output indicators & Expected & Achieved\textsuperscript{*} \\
\midrule
PhD thesis & 1 & 0 \\
Master's thesis & 2 & 1 \\
Organization of seminar or conference & 0 & 0 \\
Book & 0 & 0 \\
Book chapter & 0 & 0 \\
Conference presentation & 2 & 5 \\
Conference paper & 0 & 2 \\
Journal article & 2 & 0 \\
Other (specify): acknowledged preprint & 0 & 1 \\
\bottomrule
\end{tabular}
\par\endgroup
\textsuperscript{*}Output status updated 3 November 2026; results cutoff 30 October. PhD ongoing. One Master's thesis by Anestis Lalidis Mateo [16], HSBI, submitted 2 June 2026; completion investigator-confirmed. Five delivered presentations (ASSC 2024/2025, SPR 2025, MuC 2025, AlpCHI 2026), investigator-confirmed 3 November. Conference papers [12,13] and preprint [4] acknowledge BIAL; preprint counted only under Other. Journal acceptance unconfirmed. Software/data excluded without acknowledgement verification. Expected counts: application p.~15 [1].


\newpage
\subsubsection*{Publication list and related manuscripts}
\textbf{Abstracts (proceedings and meetings).} None verified.

\textbf{Full papers.}

Fejer G, Holzapfel T, G\'omez-Emilsson A, Hirvonen T, Jermaks R, Blum J, Lobser D, Lalidis Mateo A, Glowacki D, Gaebler M, Lenggenhager B. 2025. \emph{Viscereality: A Bio-responsive VR System for Breath-Based Interactions and Coupled Oscillator Dynamics to Augment Altered States of Consciousness}. Mensch und Computer 2025 Workshopband, MCI-WS11, 8 pp. \href{https://doi.org/10.18420/muc2025-mci-ws11-174}{doi:10.18420/muc2025-mci-ws11-174}. Bonn, Germany. Gesellschaft f\"ur Informatik e.V. \href{https://www.hsbi.de/publikationsserver/download/6574/6575/muc-2025-mci-174.pdf}{Open-access PDF}.

Fejer G, Holzapfel T, Hirvonen T, Lalidis Mateo A, Blum J, Gaebler M, Lenggenhager B. 2026. \emph{Breathing Space: Spatial Mapping of Breath and Cardiac Biofeedback for Affective State Representation and Coherence Training in Viscereality}. Proceedings of the 1st International Conference on Human-Computer Interaction in the Alps, pp.~79--86. \href{https://doi.org/10.1145/3780045.3780061}{doi:10.1145/3780045.3780061}. New York, NY, USA. Association for Computing Machinery. \href{https://www.hsbi.de/publikationsserver/download/7162/7163/3780045.3780061.pdf}{Open-access PDF}.

\textbf{Contributions included in books.} None verified.


\textbf{Public preprint and journal submission.} Fejer G, Holzapfel T, Hirvonen T, Lalidis Mateo A, Blum J, Gaebler M, Lenggenhager B. 2026. \emph{Viscereality: Bioresponsive Virtual Reality for Box Breathing and Altered-State Experience}. Research Square, v1. \href{https://doi.org/10.21203/rs.3.rs-10842634/v1}{doi:10.21203/rs.3.rs-10842634/v1}. BIAL grant 308/22 acknowledged. Under review at \emph{Virtual Reality}. \href{https://github.com/GeorgeFejer91/Viscereality-Bioresponsive-Virtual-Reality-for-Box-Breathing-and-Altered-State-Experience}{GitHub data and analysis code}. \href{https://doi.org/10.5281/zenodo.22073981}{Archived reproducibility package}.

\textbf{Manuscripts in preparation and study materials.} These entries are outside the achieved publication counts.

\emph{Respiratory Modulation of Peripersonal Space Representation}. Draft for author review. \href{https://github.com/GeorgeFejer91/Respiratory-Modulation-of-Peripersonal-Space-Representation}{GitHub manuscript and data}. \href{https://georgefejer91.github.io/Respiratory-Modulation-of-Peripersonal-Space-Representation/}{Analysis results}.

\emph{PPS Toolkit: A suite for designing, running, and recreating audio-tactile peripersonal-space experiments}. Methods manuscript in preparation for \emph{Behavior Research Methods}. \href{https://github.com/GeorgeFejer91/pps-kit}{GitHub software}. \href{https://ppskit.qzz.io/documentation/}{Documentation}.

\emph{eMotion}. Early study manuscript outline. Eight participants completed both mappings. \href{https://github.com/GeorgeFejer91/emotion-vr-data}{GitHub data and codebook}.


Delivery of the planned documentary and exhibition remains unverified.

\subsection{Journal status and explanation of delay}
\label{sec:journal}
The preprint appeared on 11 September 2026 [4]. As of 30 October, the team was still awaiting a decision from Springer's \emph{Virtual Reality} (Appendix A). The investigators report that finding a second reviewer took about six weeks. BIAL requires at least one accepted peer-reviewed journal article. Springer identifies reviewer availability as one influence on turnaround [14]. A 2026 editorial describes additional review work associated with submissions produced with generative artificial intelligence [15]. The cause of delay for this manuscript is unknown.

\subsection{Recommendations}
Verify participant-condition linkage and respiratory timing in Respiratory PPS. Before interpreting a boundary shift, establish that the task can locate a transition within the tested distance range. In eMotion, continue recruitment toward the planned 24 and verify the quality of the heart recordings for each condition before analysing cardiac responses.

\textbf{Acknowledgement.} BIAL Foundation grant 308/22. Subsequent development also received Baden-W\"urttemberg Stiftung support through Forschungswerkstatt [4].
